\documentclass[a4paper, 12pt]{article}
\usepackage{lib}

\title{Majoration du nombre de mouvement nécessaire à la résolution du jeu \url{loopover.xyz}}
\author{Parent Quentin}
\date{}

\begin{document}
	\maketitle
	
	\section*{Introduction}
	caser ici l'abstract + la problématique
	
	\section{Jeu}
		\subsection{Définition}
			objectif, définition des mouvements, notation
		\subsection{Groupe de représentation du jeu} \label{gr}
			sg de $\perm_{N^2}$, engendré par blabla
			
			À l'aide de commutateurs, on engendre les 3c $\rightarrow$ contient $\alt_{N^2}$
			
			\begin{itemize}
				\item si $N$ est pair, puisque \texttt{U1} est de signature $-1$, le groupe du puzzle est donc $\perm_{N^2}$
				\item si $N$ est impair, toues les mouvements élémentaires sont de signature 1, le groupe du puzzle est donc $\alt_{N^2}$
			\end{itemize}
	
	\section{Résolution théorique}
	\subsection{Minoration}
		Si $k$ est le majorant recherché, alors on peut associer à tout état du puzzle un mélange en moins de $k$ mouvements.
		
		Or, il existe $\displaystyle \sum_{i=0}^N (4N)^i$ mélanges en moins de $k$ mouvements.
		
		De plus, en retirant les mouvements inutiles ($\mathtt{U1\ D1} = \id$, par exemple), il n'y a que (4N-1) choix à partir du second mouvement.
		
		Finalement, il existe $\displaystyle 1 + 4N\sum_{i=0}^{k-1}(4N-1)^i = 1+4N\frac{1-(4N-1)^N}{2-4N}$.
		
		On a donc $\frac{N^2!}2 \leqslant 1+4N\frac{1-(4N-1)^k}{2-4N}$ (dans le cas $N$ impair), d'où
	
	
	\subsection{Algorithme de résolution}
	\subsection{Majoration}
	
	\section{Résolution algorithmique}
	\subsection{IDA*}
	\subsubsection*{Résultats}
	\subsection{Résolution multiphase}
	\subsubsection*{Motivation}
	
	On vérifie de même qu'en \ref{gr} que le sous-groupe $\{\sigma \in G\,|\, \text{le bloc $p*q$ est résolu}\}$ est égal à $\langle R_i, U_j\rangle_{i>p,\, j>q}$, permettant donc de chercher une solution n'utilisant que ces mouvements, réduisant donc les états explorés lors de l'utilisation de l'algorithme IDA*.
	\subsubsection*{Résultats}
	
	\section{Conclusion}
\end{document}